% Source: https://artofproblemsolving.com/community/c3h1646261_need_help
% Stable identity: aops-topic-1646261-q01
Arrange the mini-triangles in $n$ horizontal rows (taking one side of the big triangle as horizontal). Row $k$ from the top contains $k$ ``up'' triangles (pointing the same way as the big triangle) and $k-1$ ``down'' triangles, alternating. In total there are $1+2+\cdots+n=\dfrac{n(n+1)}{2}$ up triangles and $\dfrac{n(n-1)}{2}$ down triangles.

\textbf{Upper bound.} Two mini-triangles sharing a side are mirror images of each other in that side, so one is up and the other is down. Hence the triangles of a chain alternate between up and down, and a chain contains at most one more up triangle than down triangles. Its length is therefore at most
\[
2\cdot\dfrac{n(n-1)}{2}+1=n^2-n+1 .
\]

\textbf{A chain of that length.} Start with the top triangle (row $1$). In row $k\ge2$, enter through the down triangle at one end of the row lying under the last triangle used in row $k-1$, and run along the row to its other end, using all $k-1$ down triangles and $k-1$ of the $k$ up triangles; this ends at the up triangle at that end of the row, which lies over the end down triangle of row $k+1$. Alternate the direction from row to row (row $2$ leftwards, row $3$ rightwards, and so on). Row $k$ contributes $2(k-1)$ triangles, so the length is
\[
1+\sum_{k=2}^{n}2(k-1)=1+n(n-1)=n^2-n+1 .
\]
The figure shows such a chain of $13$ triangles for $n=4$; the unshaded triangles are the $n-1=3$ up triangles left out.

\begin{asy}
usepackage("amsmath");
size(6cm);
defaultpen(fontsize(10pt));
fill((1.5000,2.5981)--(2.5000,2.5981)--(2.0000,3.4641)--cycle,paleblue);
fill((1.5000,2.5981)--(2.5000,2.5981)--(2.0000,1.7321)--cycle,paleblue);
fill((1.0000,1.7321)--(2.0000,1.7321)--(1.5000,2.5981)--cycle,paleblue);
fill((1.0000,1.7321)--(2.0000,1.7321)--(1.5000,0.8660)--cycle,paleblue);
fill((1.5000,0.8660)--(2.5000,0.8660)--(2.0000,1.7321)--cycle,paleblue);
fill((2.0000,1.7321)--(3.0000,1.7321)--(2.5000,0.8660)--cycle,paleblue);
fill((2.5000,0.8660)--(3.5000,0.8660)--(3.0000,1.7321)--cycle,paleblue);
fill((2.5000,0.8660)--(3.5000,0.8660)--(3.0000,0.0000)--cycle,paleblue);
fill((2.0000,0.0000)--(3.0000,0.0000)--(2.5000,0.8660)--cycle,paleblue);
fill((1.5000,0.8660)--(2.5000,0.8660)--(2.0000,0.0000)--cycle,paleblue);
fill((1.0000,0.0000)--(2.0000,0.0000)--(1.5000,0.8660)--cycle,paleblue);
fill((0.5000,0.8660)--(1.5000,0.8660)--(1.0000,0.0000)--cycle,paleblue);
fill((0.0000,0.0000)--(1.0000,0.0000)--(0.5000,0.8660)--cycle,paleblue);
draw((1.5000,2.5981)--(2.5000,2.5981),gray(0.4));
draw((1.0000,1.7321)--(3.0000,1.7321),gray(0.4));
draw((0.5000,0.8660)--(3.5000,0.8660),gray(0.4));
draw((1.0000,0.0000)--(0.5000,0.8660),gray(0.4));
draw((1.0000,0.0000)--(2.5000,2.5981),gray(0.4));
draw((2.0000,0.0000)--(1.0000,1.7321),gray(0.4));
draw((2.0000,0.0000)--(3.0000,1.7321),gray(0.4));
draw((3.0000,0.0000)--(1.5000,2.5981),gray(0.4));
draw((3.0000,0.0000)--(3.5000,0.8660),gray(0.4));
draw((0,0)--(4,0)--(2,3.4641)--cycle,linewidth(0.9));
draw((2.0000,2.7168)--(2.0000,2.4794),red+linewidth(0.8));
draw((1.8528,2.2244)--(1.6472,2.1057),red+linewidth(0.8));
draw((1.5000,1.8507)--(1.5000,1.6134),red+linewidth(0.8));
draw((1.6472,1.3584)--(1.8528,1.2397),red+linewidth(0.8));
draw((2.1472,1.2397)--(2.3528,1.3584),red+linewidth(0.8));
draw((2.6472,1.3584)--(2.8528,1.2397),red+linewidth(0.8));
draw((3.0000,0.9847)--(3.0000,0.7474),red+linewidth(0.8));
draw((2.8528,0.4924)--(2.6472,0.3737),red+linewidth(0.8));
draw((2.3528,0.3737)--(2.1472,0.4924),red+linewidth(0.8));
draw((1.8528,0.4924)--(1.6472,0.3737),red+linewidth(0.8));
draw((1.3528,0.3737)--(1.1472,0.4924),red+linewidth(0.8));
draw((0.8528,0.4924)--(0.6472,0.3737),red+linewidth(0.8));
label("$1$",(2.0000,2.8868),fontsize(8pt));
label("$2$",(2.0000,2.3094),fontsize(8pt));
label("$3$",(1.5000,2.0207),fontsize(8pt));
label("$4$",(1.5000,1.4434),fontsize(8pt));
label("$5$",(2.0000,1.1547),fontsize(8pt));
label("$6$",(2.5000,1.4434),fontsize(8pt));
label("$7$",(3.0000,1.1547),fontsize(8pt));
label("$8$",(3.0000,0.5774),fontsize(8pt));
label("$9$",(2.5000,0.2887),fontsize(8pt));
label("$10$",(2.0000,0.5774),fontsize(8pt));
label("$11$",(1.5000,0.2887),fontsize(8pt));
label("$12$",(1.0000,0.5774),fontsize(8pt));
label("$13$",(0.5000,0.2887),fontsize(8pt));
label("\mbox{}",(min(currentpicture,user=true)+max(currentpicture,user=true))/2,black); // keeps the text after the figure black when LaTeX typesets the labels inline
\end{asy}

The largest possible number of mini-triangles in a chain is $\boxed{n^2-n+1}$.
